\chapter{Evidence, promise, and limits}
\label{ch:evidence}

Imagine a funder asking a simple question after reading the literature: what,
exactly, may this project promise? The studies can support bounded assistance,
located feedback, protected comparison, and direct reader evaluation. They do
not yet show that this combination improves an expert technical document. The
review therefore follows the decisions the product must make rather than the
chronology of the field.

The search has identified sources and a reproducible protocol. After four
missing foundations were added to the test set, the latest online run recovered
fifteen of eighteen known items through broad searches and fourteen through
deliberately indirect challenge queries. The broad-recall gate passes; the
challenge gate does not. A provenance-complete RePEc/IDEAS specialist slice is
also recorded, and Crossref resolves all twenty-six DOI-bearing bibliography
rows and all fifteen DOI-bearing load-bearing evidence rows. The scholarly job
is not finished: independent screening, full-text appraisal, held-out package
benchmarks, and external review remain incomplete. The passages that follow
say what each source can bear at its present inspection depth. That boundary
narrows the claim; it does not prevent the available evidence from changing
the design.

Figure~\ref{fig:evidence-ladder} places the present literature below the
held-out feasibility and comparative evidence that Humanvoice must still earn.
Each rung supports a stronger claim, and no lower rung can substitute for the
one above it.

\hvFigureEvidenceLadder

\section{Does assistance save time or improve work?}

The positive case for generative assistance is real but heterogeneous. In a
controlled experiment, access to ChatGPT reduced the time required for some
professional writing tasks and improved external assessments of the resulting
work \citep{noy2023experimental}. A large field study of customer-support
workers found productivity gains, with larger benefits for less-experienced
workers and smaller or absent gains for the most experienced workers
\citep{brynjolfsson2025work}. A workplace experiment across firms similarly
reports time savings in email, while leaving the composition of tasks largely
unchanged \citep{dillon2025shifting}.

These results justify treating review burden as an economic outcome rather than
assuming that any fluent assistance is valuable. They do not transfer directly
to technical proposals. A proposal reader must decide whether its mechanism,
evidence, and uncertainty warrant action. That decision can sit outside the model's capability
frontier. In a field experiment with knowledge workers, assistance improved
speed and quality on tasks inside that frontier but reduced correctness on
tasks outside it, even when the recommendations looked useful
\citep{dellacqua2026jagged}.

In the pilot, authors and readers will record their time, feedback rounds, and
substantive errors together. The evaluation lead will split those observations
by writer experience and task class instead of hiding them inside one average.
That is the product hypothesis in its useful form: Humanvoice may reduce wasted
review, even when it does not make every writing task faster.

\section{Can feedback improve the next draft?}

Writing feedback is the closest literature to the proposed intervention, and
its limits are clear. In generated stories, language models produced specific,
reasonable comments but often missed the central problem or misjudged criticism
versus encouragement \citep{rashkin2025story}. A simulated-revision study treats
the link between a comment and a useful next draft as empirical
\citep{nair2024closing}. Comparisons with human resources find potential under
detailed rubrics, but agreement varies with writer proficiency
\citep{behzad2024feedback}.

The LAMP work provides a useful magnitude check. Professional writers supplied
span-level edits to generated passages across categories such as cliche,
redundancy, poor sentence structure, missing specificity, and tense
inconsistency. A span detector reached only moderate precision, and expert
agreement was itself imperfect; edited generated text nevertheless ranked above
raw generated text in blind comparisons \citep{chakrabarty2024salvage}. The
result does not say that a detector can edit safely. It says that localized
diagnosis followed by human-controlled editing is a plausible workflow, and
that the workflow must tolerate disagreement.

Humanvoice takes three choices from this evidence. A finding includes a source
span and context, so an author can disagree with it. A substantive finding is
an editorial question, not an automatic operation. Agreement between a model
and a reference label is one diagnostic; the outcomes are the resulting
revision and the reader's response.

\section{What makes technical writing comprehensible?}

When a reader cannot tell what a technical term does, the problem is usually
the relation around the term, not the term alone. That is why the relevant
scholarship is broader than generative-AI studies. Longitudinal analysis finds
that scientific texts have become harder to read over time, while specialized
terminology is associated with fewer citations in some scientific fields
\citep{plavensigray2017readability,martinez2021terminology}. These findings do
not imply that technical terms should be removed. They imply that a document
must make its actors, mechanisms, and purpose visible around the terms.

Readability is multidimensional. Coh-Metrix and related corpus work distinguish
surface features from cohesion, referential clarity, syntactic structure, and
knowledge demands \citep{graesser2004cohmetrix,crossley2022clear}. Register
varies systematically between speech and writing, disciplines, and genres
\citep{biber1988variation}. A technical proposal should therefore be compared
with an appropriate professional corpus and reader task, not optimized toward
a universal readability number.

That finding gives the pilot a concrete reader task. Ask whether the problem,
mechanism, evidence, qualification, and action can be recovered. A
sentence-length or lexical statistic can point to a passage, but it cannot
answer those questions. Humanvoice therefore protects domain terminology and
looks instead for a missing relation, an unintroduced actor, or an avoidable
burden on the reader.

\section{What happens to trust, diversity, and integrity?}

Assistance can improve an individual output while changing the distribution of
outputs. In a controlled writing experiment, instruction-tuned assistance
reduced lexical and content diversity across writers; the model-contributed
text accounted for most of the convergence \citep{padmakumar2024writing}.
Related creative-writing work reports higher individual scores alongside lower
collective diversity \citep{doshi2024generative}. Recursive training on
synthetic text supplies a mechanism for why this matters: linguistic diversity
can decline across generations \citep{guo2023curious}.

Disclosure also changes the reader's judgment. Experiments on argumentative and
creative writing report heterogeneous responses, including lower assessments in
some conditions, when assistance is disclosed \citep{li2024disclosure}. Citation
reliability is separate: language-model introductions have produced nonexistent
or incorrect citations at material rates across disciplines
\citep{mugaanyi2024citation}. A valid identifier does not establish support.

Publication guidance places responsibility in the same location as the
product. ICMJE's recommendations keep accuracy, integrity, originality,
citations, and disclosure with human authors \citep{icmje2025ai}. NIST's AI
Risk Management Framework provides a useful vocabulary for validity,
reliability, privacy, transparency, and accountability \citep{nist2023airmf}.
Humanvoice translates those principles into concrete objects: a source
snapshot, a protected-span record, a citation verification queue, and a named
human decision. It does not treat a model output as an authorship or quality
certificate.

\section{Can a model judge the document?}

Model-based evaluation is useful as a screening instrument, but it is not a
substitute for the reader whose decision matters. The LLM-as-a-judge literature
shows sensitivity to order, verbosity, and model family
\citep{zheng2023judging,dubois2024length}. Models can also favor their own
generations \citep{panickssery2024llm}. These biases can change the writing. A
judge trained to reward a familiar style can push the product
toward the very regularity it is meant to diagnose.

Interaction studies make the missing unit of analysis visible. CoAuthor and
Wordcraft show that writing with a model is an interaction in which suggestion,
selection, editing, and rejection matter, not just a comparison of final texts
\citep{lee2022coauthor,yuan2022wordcraft}. In-the-wild traces identify
recurring behaviors such as revising intent, exploring alternatives, asking
questions, adjusting style, and injecting content \citep{mysore2025paths}. A
useful evaluation record must retain enough of that path to explain why burden
or quality changed.

If the pilot uses a model judge, the evaluation lead will first calibrate it on
human labels, randomize presentation, control length, and check another model
family. The named reader still accepts or rejects the document. The protocol
does not assume that human judgment is perfectly consistent; it records the
disagreement as part of the result.

\section{Why authorship detection is not the product}

Detection research is relevant as a boundary condition. Probability-curvature
and cross-model methods demonstrate that generated text can be distinguishable
under benchmark conditions
\citep{gehrmann2019gltr,mitchell2023detectgpt,bao2024fastdetect,hans2024binoculars}.
The failure modes are more important for this proposal. Detectors have produced
high false-positive rates for non-native English writers and formal prose, and
paraphrase can defeat several detector families
\citep{liang2023detectors,weberwulff2023testing,krishna2023paraphrasing,
sadasivan2024reliably}.

Watermarks are a provenance mechanism rather than a general detector.
Provider-side watermarking can insert a statistical signal when a cooperating
model generates the text \citep{kirchenbauer2023watermark}; it does not turn
arbitrary edited prose into a reliable authorship verdict. Humanvoice therefore
records assistance at the time it occurs and keeps watermark verification, if
an institution requires it, separate from judgments about clarity and quality.

Population-level monitoring may be statistically meaningful in a narrowly
defined corpus, as in mixture estimates of AI-modified review text
\citep{liang2024monitoring}. That is a different question from whether one
document is clear or whether one author acted properly. Humanvoice excludes
per-document authorship verdicts and keeps any future population instrument
outside the revision path.

Figure~\ref{fig:detection-boundary} turns that distinction into an interface
rule. An aggregate research instrument must not acquire a single-document
button later through convenience or product pressure.

\hvFigureDetectionBoundary

\section{The evidence-to-design synthesis}

\begin{table}[H]
\centering
\small
\begin{tabularx}{\textwidth}{@{}p{0.24\textwidth}L p{0.28\textwidth}@{}}
\toprule
Evidence question & What the evidence supports & Product consequence \\
\midrule
Does assistance save work? & Effects are task- and worker-dependent; some
frontier failures reduce correctness. & Measure reader time and meaning errors,
and report heterogeneous outcomes. \\
Can feedback improve revision? & Local comments can help, but central defects
and priority are often missed; expert agreement is imperfect. & Locate findings,
show context, and leave repair and acceptance with a human. \\
What is comprehensible? & Cohesion, referential clarity, genre, and knowledge
demand matter more than one surface formula. & Use genre-aware diagnostics and
a direct purpose-and-action reader rubric. \\
What risks accompany assistance? & Diversity, disclosure, citation reliability,
and provenance can change even when prose improves. & Protect meaning-bearing
objects, verify citation identity and support, and record disclosure. \\
Can a model judge? & Judges vary with order, length, family, and self-preference. &
Use model evaluation only as calibrated screening; never as acceptance. \\
\bottomrule
\end{tabularx}
\caption{The literature changes the design rather than merely decorating it with
citations.}
\label{tab:evidence-design}
\end{table}

\section{The questions still open}

No study in the current record answers whether Humanvoice reduces feedback
rounds for expert technical documents, whether authors will tolerate its
findings, or whether the protected comparison saves more attention than it
consumes. Those are project hypotheses. The team must test them with a
comparator, a reader protocol, and a corpus with independent labels.

The current evidence status is intentionally visible but compact. Of fifteen
release gates, nine pass. The six blockers are challenge-query recall,
independent screening by two people, full-text claim anchors, design-specific
appraisal, held-out package effectiveness, and external review. Specialist
coverage and bibliographic DOI identity have passed for the bounded records
described above; neither pass turns an abstract-level record into a fully
appraised study. The proposal
therefore uses phrases such as "supports the design" and "open question" rather
than "proves" or "validated." Completing the evidence work is part of the
implementation program; it is not a reason to put search administration in
front of the product story.
